\chapter{The engine}

\label{ch:engine}

This chapter answers what a run is. It covers the five states a run can be in, what happens between
the moment a graph is accepted and the moment its terminal is durable, how the scheduler decides
which nodes may run together, what the journal records and when it is flushed, and how a run that
was interrupted is resumed rather than restarted. The engine is the part of pii that must not be
interesting: it starts a run, walks the graph, and writes down what it decided. Its complexity is
almost entirely in the two properties that follow from that --- a run must not hang, and a run must
not lose work it already did.

Both properties cost something. The no-hang property is why every arm of the tick loop returns and
why a node that can never become ready fails the run by name instead of waiting. The no-loss
property is why the journal is write-behind rather than write-through, why the terminal event is
appended before the status is flipped, and why the resume path replays the whole journal instead of
trusting a saved position.

\section{The record of a run}

A run is identified by a \texttt{RunId} and described by a \texttt{RunHeader} that is written when
the run is accepted. The header is small on purpose: it is read by every list command and by every
watchdog, so it must be one record and not a fold over the journal.

\begin{center}
\footnotesize
\begin{tabular}{@{}p{0.18\textwidth}p{0.72\textwidth}@{}}
	\toprule
	\textbf{Field} & \textbf{Meaning} \\
	\midrule
	\texttt{run} & the run id; the key every other table uses \\
	\texttt{graph} & the graph name, or the hash of its body when it has no name \\
	\texttt{status} & one of the five states in the next section \\
	\texttt{seq} & the next journal sequence number to be written \\
	\texttt{dur\_ms} & wall-clock duration; zero while running \\
	\texttt{revision} & the build revision of the binary that accepted the run \\
	\texttt{process\_start\_ms} & when that process started, so a restart is visible \\
	\texttt{created\_ms} & when the run was accepted \\
	\bottomrule
\end{tabular}
\end{center}

The last two fields exist because of an operational question that a single timestamp cannot answer:
\enquote{was this run started by the process that is running now}. Without \texttt{revision} and
\texttt{process\_start\_ms}, a run that fails after a deploy is indistinguishable from a run that
failed before it, and a bug that was fixed looks like it is still there. Rows written before the
field existed decode with an empty revision, and are counted as unattributed rather than as a
deploy.

\section{Five states}

\begin{center}
\resizebox{\linewidth}{!}{%
\begin{tikzpicture}[node distance=11mm and 16mm]
	\node[b] (running) {Running};
	\node[b, right=of running] (suspended) {Suspended\\\scriptsize carries a resume key};
	\node[b, below=of running] (succeeded) {Succeeded};
	\node[b, below=of suspended] (failed) {Failed};
	\node[b, left=of failed] (refused) {Refused\\\scriptsize refused before any spawn};

	\draw[f] (running) -- node[p,above=2pt,midway]{a node suspends} (suspended);
	\draw[bk] (suspended) -- node[p,below=2pt,midway]{resume key fires} (running);
	\draw[f] (running) -- (succeeded);
	\draw[f] (running) -- node[p,right=2pt,midway]{failure, cancel} (failed);
	\draw[sf] (running) to[bend right=25] node[p,below=2pt,midway]{agent call refused} (refused);
\end{tikzpicture}}
\captionof{figure}{The whole life cycle; there is no \emph{Pending} and no \emph{Cancelled} state:
a run exists only once it has been accepted, and a cancelled run is closed as \emph{Failed} with
\texttt{cancelled:} in its error text, which is what the record has always meant.}
\end{center}

\begin{itemize}
	\item \textbf{Running} --- executing, or waiting to be resumed by a caller that already holds a
	key. It is also the state a cancelled run is left in if the process dies mid-cancel, which is
	deliberate: recovery finds it.
	\item \textbf{Suspended} --- the run parked itself on a resume key. A timer, for instance,
	never blocks a thread; it journals a \texttt{Suspended} event and returns the key the scheduler
	will fire.
	\item \textbf{Succeeded} --- the run reached a terminal disposition with a value.
	\item \textbf{Failed} --- the run ended in an error: a node failed terminally, a limit was
	exceeded, a join could never be satisfied, or an operator cancelled it.
	\item \textbf{Refused} --- an agent call was refused before any process was spawned. It is
	distinct from \emph{Failed} because nothing was attempted: a refusal is a decision, and the
	store records it as one.
\end{itemize}

The absence of a \emph{Cancelled} status is not an oversight. Cancellation has two forms, and
neither needs a sixth state. A run that is not executing is closed as \emph{Failed} with
\texttt{cancelled: <reason>} in its error, which is what an operator reading the record needs to
know. A run whose node \emph{is} executing is refused by default --- the daemon answers
\texttt{run.cancel\_refused} and names the node --- because cancelling would abandon work
half-done. The explicit kill form accepts that, waits up to thirty seconds for the node to leave
flight, and still leaves the run open if it has not: a run is never closed while its agent may
still be running.

\section{Starting a run}

\texttt{spawn\_run} returns in under fifty milliseconds, and everything that must be true before
the caller gets an id is done on the calling thread. The order matters, because each step can refuse
and the refusal has to leave a coherent record.

\begin{enumerate}
	\item The runtime is checked. A start from a plain thread is refused by name rather than
	panicking on a run thread, which would crash-loop a supervisor.
	\item The run is registered as queued, with its own cancel token, so a cancel that arrives
	before the pool picks it up is not lost.
	\item An optional \texttt{scope\_root} in the input confines the run's write scope to one
	directory. An invalid root is rejected here, before anything runs, rather than producing a
	permissive run.
	\item The context is built with the graph's \texttt{[limits]} as a real budget --- tokens, tool
	calls, micro-dollars and wall clock --- and with the capability set the run's scope allows.
	\item The header and the \texttt{Started} event are written and flushed, so the id the caller
	receives is already durable.
	\item A quarantine check may refuse the run; the refusal is recorded as a deny decision before
	the error is returned.
	\item The work is submitted to the run pool. The closure checks whether the run was cancelled
	while queued and returns without executing if so.
\end{enumerate}

The fifth step is the one with a visible consequence: because the header is durable before the
caller returns, a run id can be used immediately --- to subscribe, to cancel, to read the trace ---
without a race against the pool.

\begin{lstlisting}[language=bash]
pii run start queue --input '{"workdir":"/projects/pi-improved","base":"<40-hex>"}'
pii run ls --status Suspended          # what is waiting, and on what
pii run get <run_id>                   # header: graph, status, duration
pii run trace <run_id>                 # the journal, in sequence order
\end{lstlisting}

\section{The tick loop}

The loop is one function, and its shape is the answer to \enquote{how does a graph make progress}.
Each iteration either dispatches a batch, or ends the run. There is no arm that waits.

\begin{center}
\resizebox{\linewidth}{!}{%
\begin{tikzpicture}[node distance=9mm and 14mm]
	\node[b] (top) {tick};
	\node[v, below=of top] (cancel) {cancelled?};
	\node[v, below=of cancel] (limits) {limits ok?};
	\node[v, below=of limits] (join) {unsatisfiable\\\scriptsize join?};
	\node[b, below=of join] (batch) {drain every ready node\\\scriptsize one batch};
	\node[v, below=of batch] (empty) {batch empty?};
	\node[b, right=of batch] (part) {pure chunk, then\\\scriptsize serial, applying as results arrive};
	\node[b, left=of empty] (term) {terminal};

	\draw[f] (top) -- (cancel);
	\draw[f] (cancel) -- node[p,right=2pt,midway]{no} (limits);
	\draw[f] (limits) -- node[p,right=2pt,midway]{yes, dead to a fixpoint} (join);
	\draw[f] (join) -- node[p,right=2pt,midway]{none} (batch);
	\draw[f] (batch) -- (empty);
	\draw[f] (empty) -- node[p,above=2pt,midway]{no} (part);
	\draw[sf] (part.south) to[bend left=30] node[p,below=2pt,midway]{next tick} (top.west);
	\draw[f] (empty) -- node[p,above=2pt,midway]{yes} (term);
	\draw[bk] (cancel.west) to[bend right=20] node[p,left=2pt,midway]{flush, stop} (term.north);
	\draw[bk] (join.east) to[bend left=20] node[p,right=2pt,midway]{yes} (term.north east);
\end{tikzpicture}}
\captionof{figure}{One iteration; the two red exits are the no-hang guarantee: a cancellation and
an unsatisfiable join both end the run inside the same tick, with a recorded reason.}
\end{center}

In order, one iteration does this:

\begin{enumerate}
	\item \textbf{Cancellation.} If the run's token is set, the journal is flushed and the run stops.
	The flush is not optional: without it a cancel before the write-behind threshold would lose every
	completed node, and a resume would re-run work whose side effects already happened.
	\item \textbf{Limits.} The step count and the wall clock are checked first, then the three
	consumable budgets. A budget exhausted by any node is visible here on the next tick, because
	nodes share the run's budget rather than holding their own.
	\item \textbf{Dead propagation.} A node all of whose inbound data edges are dead is itself dead,
	and so on to a fixpoint. Each newly dead node journals a \texttt{NodeDead} event, so a branch
	that was skipped is distinguishable from a branch that never existed.
	\item \textbf{Unsatisfiable join.} A join that can never be satisfied fails the run immediately
	and names the node. This is checked every tick rather than at the end, because a quorum that has
	already lost too many arms cannot recover.
	\item \textbf{The batch.} Every node that is ready and not already executed, dead, failed or
	cancelled is drained into one batch. The batch is journaled as a \texttt{TickStarted} event with
	its width, which is what makes a wide fan-out visible in a trace.
	\item \textbf{Budget and entry checks at dispatch.} A node whose \texttt{loop\_budget} is
	already spent, or a call head entered more than \texttt{max\_entries} times, fails the run here
	--- before the batch runs, not after.
	\item \textbf{Inputs frozen.} Each node's input is built from the current edge state and the
	blackboard, once, before any node in the batch runs. A node therefore cannot observe a sibling's
	output from the same tick.
	\item \textbf{Partition.} Nodes whose effects are parallel-safe run in bounded chunks; nodes that
	write the filesystem or mutate an external system run one at a time, after the parallel batch.
	\item \textbf{Apply as results arrive.} Each result is applied to the scheduling state
	immediately, and race losers are cancelled after \emph{each} result rather than at the end of the
	batch, so a losing branch does not keep burning a timer.
	\item \textbf{Restart or end.} A terminal failure recorded mid-batch ends the run now; a
	suspension ends it as \emph{Suspended}; otherwise the loop iterates.
\end{enumerate}

The choice to drain \emph{all} ready nodes rather than one per tick is what makes the loop's cost
proportional to depth rather than to nodes. It also means a graph's fan-out is a real budget
question: a hundred pure nodes ready at once will use up to \texttt{max\_parallel} threads, and the
single-element chunk case is dispatched inline precisely because spawning one OS thread per node
showed up as retained memory under soak load.

\section{Readiness, and the four joins}

A node is ready when its inbound edges allow it. The rules are ordered, and the order is the part
that is easy to get wrong.

\begin{enumerate}
	\item The entry node is ready on its first pass even if it has inbound edges. A loop's back
	edge is pending at the start, and without this exception a cyclic graph would deadlock on tick
	one.
	\item Any pending \emph{order} edge makes a node not ready. Order edges gate readiness before
	data is considered at all.
	\item Zero inbound data edges means ready. This must be tested before the join arms, because
	\enquote{all of nothing} is true and would otherwise report the entry as dead.
	\item Otherwise the node's join policy decides.
\end{enumerate}

\begin{center}
\footnotesize
\begin{tabular}{@{}p{0.16\textwidth}p{0.36\textwidth}p{0.36\textwidth}@{}}
	\toprule
	\textbf{Join} & \textbf{Ready when} & \textbf{Never runs when} \\
	\midrule
	\texttt{all} (default) & every inbound data edge fired & every edge dead (dead node), or any edge
	dead while others fired (unsatisfiable) \\
	\texttt{any} & at least one edge fired & every edge dead \\
	\texttt{any:cancel} & as \texttt{any}, and the losers are cancelled once one wins & every edge
	dead \\
	\texttt{quorum:K} & at least $K$ edges fired & fewer than $K$ can still fire \\
	\texttt{allsettled} & every edge settled, fired or dead & every edge dead \\
	\bottomrule
\end{tabular}
\end{center}

The difference between \texttt{all} and \texttt{allsettled} is the difference between \enquote{all
my inputs succeeded} and \enquote{all my inputs are decided}. A node that must proceed whether or
not one branch died uses \texttt{allsettled}; a node that needs every value uses \texttt{all}. The
static validator refuses an \texttt{any} or \texttt{quorum} join with fewer than two inbound data
edges --- a one-armed race is a typo --- but deliberately does not over-enforce a single-armed
\texttt{allsettled}, which is a legitimate idiom.

\begin{center}
\resizebox{\linewidth}{!}{%
\begin{tikzpicture}[node distance=8mm and 12mm]
	\node[b] (a) {branch A};
	\node[b, below=of a] (b) {branch B};
	\node[m, right=14mm of a] (join) {join};
	\node[b, right=of join] (next) {next node};
	\node[v, below=of next] (dead) {dead branch};

	\draw[f] (a.east) -- node[p,above=2pt,midway]{ok} (join.north west);
	\draw[f] (b.east) -- node[p,below=2pt,midway]{done} (join.south west);
	\draw[f] (join) -- (next);
	\draw[bk] (b.south) -- (dead.north);
	\draw[sf] (dead.west) to[bend left=14] node[p,below=2pt,midway]{all dead?} (join.south);
\end{tikzpicture}}
\captionof{figure}{Why the join policy is part of the graph rather than a scheduler default; two
branches may both feed one node, and whether that node waits for both, for one, or for both to be
\emph{decided} is a fact about the graph that no scheduler can infer.}
\end{center}

\section{The journal}

The journal is the primary evidence, and it is written by one component with one sequence counter.
Every event is appended with a sequence number computed from the header and the pending buffer, the
observer is notified immediately, and the event is buffered for a durable write in one transaction
with the header.

\begin{longtable}{@{}p{0.26\textwidth}p{0.66\textwidth}@{}}
	\toprule
	\textbf{Event} & \textbf{What it records} \\
	\midrule
	\endhead
	\texttt{Started} & the graph name and the hash of the run input \\
	\texttt{TickStarted} & the batch of nodes drained in one tick, and its width \\
	\texttt{NodeStarted} & a node began, at a millisecond offset \\
	\texttt{NodeFinished} & the node, the port it took, the value, the value hash, offset and
	duration \\
	\texttt{NodeFailed} & the error, whether it is retryable, and the attempt number \\
	\texttt{NodeProgress} & a heartbeat: elapsed time and the child's CPU, so a wedged node (time
	grows, CPU flat) reads differently from a working one \\
	\texttt{NodeCancelled}, \texttt{NodeDead} & a race loser, and a node killed by dead-branch
	propagation \\
	\texttt{NodeInterrupted} & an audit marker written on resume for a node the previous attempt
	left mid-flight; it carries no state transition \\
	\texttt{Suspended}, \texttt{Resumed} & the resume key, and the hash of the payload that
	delivered it \\
	\texttt{NodeObserved} & an event a node's agent emitted, recorded from the agent's own events
	rather than from the node's return value \\
	\texttt{EffectStarted}, \texttt{EffectConfirmed} & an idempotency key, written \emph{before} an
	effect and confirmed after it \\
	\texttt{NodeRetried} & an operator re-ran one node; earlier terminal events are history \\
	\texttt{RunSpawned} & a \texttt{spawn} node launched a child run \\
	\texttt{Finished} & the terminal and the hash of its output \\
	\bottomrule
\end{longtable}

Two of these groups carry more weight than the rest. The \texttt{EffectStarted} /
\texttt{EffectConfirmed} pair is the only defence against an effect that happened but was not
recorded: a key that was started and never confirmed is reported as \texttt{effect.unknown}, and a
resume refuses the run rather than re-performing an effect whose outcome nobody knows. And
\texttt{NodeObserved} is how the journal learns what an agent did without trusting the node's return
value: the fields are the agent's event fields, copied, not paraphrased.

\subsection*{Write-behind, and the two moments it is flushed}

Journaling is buffered. The writer keeps a pending list, flushes when the list reaches its batch
threshold or when the sequence number crosses a snapshot boundary every thirty-two events, and
flushes unconditionally at three points that matter:

\begin{itemize}
	\item before a terminal is recorded --- the \texttt{Finished} event is appended \emph{before} the
	status is flipped, so a reader that polls the header cannot see a finished run while the terminal
	event is still in the buffer;
	\item when a run is cancelled, so completed work survives;
	\item on resume, before anything new is appended, so the replayed prefix is the durable one.
\end{itemize}

\begin{center}
\resizebox{\linewidth}{!}{%
\begin{tikzpicture}[node distance=9mm and 13mm]
	\node[b] (a1) {append};
	\node[b, right=of a1] (a2) {append};
	\node[b, right=of a2] (a3) {append\\\scriptsize 32nd};
	\node[b, right=of a3] (flush) {flush\\\scriptsize one transaction};
	\node[b, below=of a2] (obs) {observer\\\scriptsize every append};
	\node[b, below=of flush] (finish) {\texttt{Finished}\\\scriptsize appended first};
	\node[b, right=of finish] (status) {status flipped};
	\node[b, below=of status] (durable) {durable\\\scriptsize both, one flush};

	\draw[f] (a1) -- (a2);
	\draw[f] (a2) -- (a3);
	\draw[f] (a3) -- (flush);
	\draw[sf] (a2) -- (obs);
	\draw[f] (finish) -- (status);
	\draw[f] (status) -- (durable);
	\draw[sf] (flush.south) to[bend right=10] (finish.north);
\end{tikzpicture}}
\captionof{figure}{Two clocks in the writer: the live view advances on every append, the durable view on a flush, and
the terminal path appends before flipping the status so that both land in one transaction.}

\end{center}

The failure mode this design accepts is bounded and stated: a process that dies between two flushes
loses the buffered events, and the nodes they described will run again on resume. That is why
effects are journalled separately and before they happen --- the one category where \enquote{run
again} is not acceptable gets its own protocol.

\section{How a run ends}

When the batch is empty, the engine decides the terminal. Three rules, in order:

\begin{enumerate}
	\item If a failure was recorded during dispatch, it is the terminal. Nothing overrides it.
	\item If any node is neither executed, dead, failed nor cancelled and is not ready, the run fails
	with \texttt{no progress: node '<id>' never became ready}. This is the message an operator sees
	when an edge is missing, and it names the node rather than the edge, because the edge is usually
	the thing that is not there.
	\item Otherwise the output is the value of the latest finished node that has no outbound data
	edge on the port it took --- the value nobody consumed. If several qualify, the latest by
	sequence wins.
\end{enumerate}

The third rule has a consequence that surprises graph authors: a node that leaves on a
non-\texttt{ok} port carrying an error makes the run a failure even if the graph \enquote{reached
the end}. The check
is deliberate. A success terminal that carries an error is the shape of a bug that reports green
while the work failed, and it is refused at the point where the terminal is built rather than left
for a reader to notice.

\section{Suspension and resume}

A node that must wait --- a timer, an approval, a peer message --- does not block. It returns a
suspension: a reason and a key. The run is finished as \emph{Suspended}, and something outside the
engine owns the key: the timer wheel fires it, an operator resolves the approval, a peer message
arrives.

Resuming is not \enquote{continue from where we were}, because there is no \enquote{where we were}
after a restart. It is a replay:

\begin{enumerate}
	\item The journal is read in sequence order and the scheduling state is rebuilt: which edges
	fired, which nodes executed, which are dead.
	\item A node with a \texttt{NodeFinished} is never re-invoked; its value is in the journal and
	goes back on the blackboard.
	\item A node with a \texttt{NodeStarted} and no terminal event \emph{is} re-run. This is the
	documented boundary: the engine knows the node began and does not know whether its side effect
	happened, so it runs it again unless the effect protocol says otherwise.
	\item Before any of that, an \texttt{EffectStarted} with no matching \texttt{EffectConfirmed}
	refuses the resume with \texttt{effect.unknown}.
	\item A confirmed effect whose \texttt{NodeFinished} never landed is completed from the journal
	rather than re-performed.
	\item The nodes the journal shows as still suspended are dispatched through the resume entry
	point rather than a fresh run, so a fired timer resumes to \enquote{next} instead of
	suspending again.
\end{enumerate}

\begin{center}
\resizebox{\linewidth}{!}{%
\begin{tikzpicture}[node distance=8mm and 12mm]
	\node[s] (journal) {journal};
	\node[b, right=of journal] (replay) {replay\\\scriptsize rebuild edge states};
	\node[v, right=of replay] (effect) {effect\\\scriptsize unknown?};
	\node[b, right=of effect] (refuse) {refuse the resume};
	\node[b, below=of effect] (resume) {dispatch the\\\scriptsize suspended nodes};
	\node[b, right=of resume] (continue) {tick loop\\\scriptsize continues};
	\node[b, below=of replay] (rerun) {re-run started\\\scriptsize without a terminal};

	\draw[f] (journal) -- (replay);
	\draw[f] (replay) -- (effect);
	\draw[bk] (effect) -- node[p,above=2pt,midway]{yes} (refuse);
	\draw[f] (effect) -- node[p,right=2pt,midway]{no} (resume);
	\draw[f] (resume) -- (continue);
	\draw[sf] (replay) -- (rerun);
\end{tikzpicture}}
\captionof{figure}{Resume is a replay with two exceptions --- finished nodes are not re-run, and suspended ones get
their own entry point --- so everything else is reconstructed from events.}

\end{center}

\section{Recovery at boot}

A restart is the normal case, not the exceptional one. Three mechanisms cover it, and each has a
distinct failure it was built for. The operator's side of all three is small:

\begin{lstlisting}[language=bash]
pii run resume <run_id>                 # replay the journal and continue
pii run retry <run_id> --node <id>      # re-run one node; its earlier events stay as history
pii run abandon <run_id> --reason '...' # close a run nothing can finish, with a reason on record
pii run cancel <run_id> [--kill]        # refused while a node executes, unless killed
\end{lstlisting}

\paragraph{Agents left running.} A daemon that dies in the middle of an agent leaves the agent with
no owner. Before the runs it left \emph{Running} are closed, boot recovery kills each agent it
started: a container by its \texttt{io.pii} and \texttt{pii.run} labels, a host process by the pid
journalled while \texttt{/proc} still shows that pid running the same program. Every reap is
journalled on its run as \texttt{coder.orphan\_reaped}, and the run's failure text says an agent was
killed --- so the operator sees why the run ended, rather than an unexplained failure.

\paragraph{Suspended runs with no timer.} A timer row used to be deleted before the resume was
spawned, so a process that exited in between left a run \emph{Suspended} with nothing that would
ever fire. The fixes are small and together they close the hole: the row is deleted only after the
resume is spawned, a failed spawn puts the row back or fails the run with
\texttt{run.resume.failed}, and the daemon calls the existing recovery pass once at boot rather than
only after a lagging bus
event. A run must never be left suspended with no timer and no event saying so.

\paragraph{Queue claims.} A row claimed by a run that then died used to block the queue for ever,
because a \texttt{Running} row makes \texttt{claim\_next} refuse to hand out anything. The claim is
now a lease: the item records who claimed it and when, and it is requeued to \texttt{pending} with a
\texttt{requeued\_after\_restart} event when the claiming run is terminal or no longer exists. It is
\emph{not} requeued while its run is alive or suspended, which is the near-miss control the fix
needed: a suspended run still blocks its row, because its timer will resume it.

\begin{center}
\resizebox{\linewidth}{!}{%
\begin{tikzpicture}[node distance=8mm and 12mm]
	\node[b] (boot) {daemon starts};
	\node[b, right=of boot] (reap) {reap agents\\\scriptsize labels, or pid};
	\node[b, right=of reap] (suspend) {re-arm suspended\\\scriptsize runs without a timer};
	\node[b, right=of suspend] (queue) {requeue claims\\\scriptsize whose run is gone};
	\node[s, below=of reap] (j1) {journal};
	\node[s, below=of suspend] (j2) {timer wheel};
	\node[s, below=of queue] (j3) {queue store};

	\draw[f] (boot) -- (reap);
	\draw[f] (reap) -- (suspend);
	\draw[f] (suspend) -- (queue);
	\draw[f] (reap) -- (j1);
	\draw[f] (suspend) -- (j2);
	\draw[f] (queue) -- (j3);
\end{tikzpicture}}
\captionof{figure}{Boot, in order: each step writes to the store that owns the fact, and nothing is fixed in memory
alone.}

\end{center}

\section{Failure signatures}

The table is the practical part of this chapter: what a failure looks like in the record, and what
it usually means. Every signature is a real string or event in the code.

\begin{longtable}{@{}p{0.34\textwidth}p{0.30\textwidth}p{0.28\textwidth}@{}}
	\toprule
	\textbf{What you see} & \textbf{Where} & \textbf{Usually means} \\
	\midrule
	\endhead
	\texttt{no progress: node 'X' never became ready} & terminal error, status
	\emph{Failed} & an inbound edge that can never fire; a branch that dead-ended \\
	\texttt{join can never be satisfied: node 'X'} & terminal, within one tick & two arms on mutually
	exclusive ports of one router \\
	\texttt{max\_steps exceeded} & terminal, status \emph{Failed} & a cycle without a
	\texttt{loop\_budget} that validation did not catch, or a node that keeps re-entering \\
	\texttt{max\_wall\_ms exceeded} & terminal & a blocking node with no timeout of its own \\
	\texttt{loop\_budget exceeded for node 'X'} & terminal, mid-batch & the budget is exactly the
	number of allowed runs; the failure is the run after it \\
	\texttt{call 'X' entered more than max\_entries} & terminal, before dispatch & an inner call
	inside an outer loop, whose count is cumulative for the run \\
	\texttt{effect.unknown} & refusal on resume & an effect began and was never confirmed \\
	no \texttt{NodeFinished} after a \texttt{NodeStarted} & journal tail & the process died
	mid-node; the node will re-run \\
	a run \emph{Running} with a dead process & header, plus reaps & crash; boot recovery closes it
	after killing the agent \\
	\bottomrule
\end{longtable}

Two of these deserve a note. \texttt{no progress} is the most common message in this repository and
almost always means an edge is missing rather than that the scheduler is wrong; the node it names is
the one that could not proceed, not the one that failed to feed it. And \texttt{loop\_budget
exceeded} fires on the run \emph{after} the allowed number, so a budget of three produces exactly
three executions and then a failure --- which is the behaviour a bounded loop wants, and the reason
the check is at dispatch rather than at re-entry.

\section{The limits, as a table}

\begin{center}
\footnotesize
\begin{tabular}{@{}p{0.20\textwidth}p{0.22\textwidth}p{0.48\textwidth}@{}}
	\toprule
	\textbf{Key} & \textbf{Default} & \textbf{What it bounds} \\
	\midrule
	\texttt{max\_steps} & 1000 & scheduler iterations, not node runs; a wide graph uses few \\
	\texttt{max\_wall\_ms} & 600\,000 & elapsed time from acceptance \\
	\texttt{max\_tokens} & unbounded & tokens charged by agent nodes, shared by the run \\
	\texttt{max\_tool\_calls} & unbounded & tool calls across the run; exhaustion routes the agent
	to its \texttt{budget} port \\
	\texttt{max\_micro\_usd} & unbounded & spend, in micro-dollars \\
	\texttt{max\_parallel} & 0 & parallel-safe nodes per tick; zero means the host's available
	parallelism \\
	\bottomrule
\end{tabular}
\end{center}

The defaults are chosen so that an unattended run cannot loop for ever or spend without a ceiling,
while a graph that says nothing still behaves like a normal program. A graph that needs longer says
so in its own file, where the number is reviewable --- which is the whole reason the limits are part
of the graph rather than a daemon-wide setting.
